\section{Methodology}
\subsection{Part 1: Policy Gradients on Hopper}
The first part uses \texttt{Hopper-v4}, a continuous-control MuJoCo task. The action is continuous and the policy is stochastic. In \texttt{part1/agent.py}, the policy network is a two-hidden-layer MLP with 64 units per layer and \texttt{tanh} activations. The actor outputs the mean of a diagonal Gaussian distribution; the standard deviation is a learned parameter passed through \texttt{softplus}. All Part 1 experiments use Adam with learning rate $3 \times 10^{-4}$ and discount factor $\gamma=0.99$.

\paragraph{REINFORCE.}
For an episode trajectory, the implementation computes Monte-Carlo returns
\begin{equation}
G_t = \sum_{k=t}^{T} \gamma^{k-t} r_k ,
\end{equation}
and minimizes the negative policy-gradient objective
\begin{equation}
\mathcal{L}_{\mathrm{PG}} = - \sum_t (G_t-b)\log \pi_\theta(a_t|s_t),
\end{equation}
where $b$ is either zero or a constant baseline. 

\paragraph{Actor-Critic Variant.}
The actor-critic implementation uses the same actor network plus a critic network that predicts $V_\phi(s)$. At each update, the critic is trained by regression to the one-step bootstrapped target
\begin{equation}
y_t = r_t + \gamma V_\phi(s_{t+1})(1-d_t),
\end{equation}
where $d_t$ indicates episode termination and the next-state value is detached from the gradient. The actor uses the corresponding advantage estimate
\begin{equation}
A_t = y_t - V_\phi(s_t),
\end{equation}
which is normalized over the collected episode before computing the policy-gradient loss. 

\subsection{Part 2: PandaPush Pipeline}
The second part uses the \texttt{panda-gym} environment. The source object mass is $1.0$ kg, while evaluation is performed on two target settings described in the evaluation protocol. The goal position is fixed at the center of the table, while the object initial position is randomized in the $xy$ plane. The training scripts use the dense task reward and an episode horizon of 500 steps.

\paragraph{Algorithms.}
For Part 2 we rely on the Stable-Baselines3 implementations of two state-of-the-art actor-critic algorithms. PPO~\cite{schulman2017ppo} is an \textbf{on-policy} method that collects rollouts under the current policy and updates it by maximizing a clipped surrogate objective,
\begin{equation}
\mathcal{L}^{\mathrm{CLIP}}(\theta) = \mathbb{E}_t\!\left[\min\!\big(\rho_t A_t,\ \mathrm{clip}(\rho_t, 1-\epsilon, 1+\epsilon)\, A_t\big)\right],
\end{equation}
where $\rho_t = \pi_\theta(a_t|s_t)/\pi_{\theta_{\mathrm{old}}}(a_t|s_t)$ is the probability ratio, $A_t$ is a GAE advantage estimate, and the clip range $\epsilon$ bounds the policy update. SAC~\cite{haarnoja2018sac} is an \textbf{off-policy} maximum-entropy method that reuses past transitions from a replay buffer and optimizes the entropy-regularized objective
\begin{equation}
J(\pi) = \mathbb{E}\!\left[\sum_t r_t + \alpha\, \mathcal{H}\big(\pi(\cdot|s_t)\big)\right],
\end{equation}
where $\mathcal{H}$ is the policy entropy and the temperature $\alpha$ (tuned automatically) trades off reward and exploration. Because the manipulation observation is a dictionary (achieved goal, desired goal, and state), both algorithms use the \texttt{MultiInputPolicy}.

\paragraph{Training Setting}

The training was performed trying different hyperparameters, and then choosing 
the ones that best performed. Table~\ref{tab:training-hyperparameters} reports the main hyperparameters used for the final SAC and PPO runs.

\begin{table}[H]
\centering
\small
\resizebox{\linewidth}{!}{%
\begin{tabular}{lcc}
\toprule
Hyperparameter & SAC & PPO \\
\midrule
Training steps & 500k & 500k \\
Policy & \texttt{MultiInputPolicy} & \texttt{MultiInputPolicy} \\
Learning rate & $10^{-3}$, $3e^{-4}$, $10^{-4}$ & $10^{-3}$, $3e^{-4}$, $10^{-4}$ \\
Discount $\gamma$ & 0.95, 0.98, 0.99 & 0.95, 0.98, 0.99 \\
Batch size & 512, 1024, 2048 & 512, 1024, 2048 \\
Gradient steps / epochs & -1 & 5 \\
Rollout steps & -- & 2048 \\
Learning starts & 100k & -- \\
Replay buffer size & $10^6$ & -- \\
Target update $\tau$ & 0.05 & -- \\
Entropy coefficient & auto & 0.01 \\
GAE $\lambda$ & -- & 0.95 \\
Clip range & -- & 0.2 \\
Policy network & default SB3 & [512, 256, 128] \\
\bottomrule
\end{tabular}
}
\caption{Main hyperparameters used for the Part 2 Stable-Baselines3 training runs. Dashes indicate algorithm-specific parameters that are not used.}
\label{tab:training-hyperparameters}
\end{table}

\paragraph{Reward Shaping.}

The wrapper \texttt{part2/wrappers.py} changes only the scalar training reward, not the dynamics. The shaped reward is
\begin{equation}
r_{\mathrm{train}} = r_{\mathrm{env}} - p + B \mathbf{1}[d(o,g)<\epsilon],
\end{equation}
where:
\begin{itemize}
    \item $p$ is a per-step penalty,
    \item $B$ is a close-goal bonus,
    \item $d(o,g)$ is the object-goal distance
\end{itemize}  
The values used for the training are $p=10^{-2}$, $B=1.0$, and $\epsilon=0.05$. 

Reward shaping was added on top of the dense environment reward to encourage faster goal reaching through a small per-step penalty and an additional bonus when the object is close to the target.
Evaluation uses the environment reward directly, so success rate is not computed from the shaped reward.

\paragraph{Domain Randomization.}
The wrapper \texttt{part2/rand\_wrapper.py} implements three modes. In \texttt{none}, the nominal environment dynamics are kept. In \texttt{udr}, the wrapper samples a new mass and friction values at every reset. The current UDR training ranges are:
\begin{equation}
m \sim \mathcal{U}(1,5), \quad
\mu_{\mathrm{obj}},\mu_{\mathrm{table}} \sim \mathcal{U}(0.5,1.2),
\end{equation}
and object spinning friction is sampled from $\mathcal{U}(0,0.005)$.
In \texttt{adr}, the wrapper starts from the nominal mass and friction and adapts the sampling ranges using a window of 20 episode outcomes. If the mean success rate is at least 0.70, the mass upper bound expands toward 5 kg and the friction ranges widen toward their global limits. If the mean success rate is at most 0.35, the ranges contract back toward the nominal values. With the default step size, the mass range expands by 1 kg per successful update, while lateral-friction ranges expand by 0.175 per update.

\paragraph{Evaluation Protocol.}
Each model is evaluated with deterministic actions for 50 episodes, using seeds 0, 42, and 99. The main metric is the final success rate, computed from \texttt{info["is\_success"]} at the end of each episode. We use two target settings: an easy fixed target with the nominal target mass, and a hard target where mass and friction are randomized at test time,
from a uniform distribution. We took this direction in order to test the model in various and unknow settings. This was done to achieve a single policy capable of generalizing 
different materials (both for the surface and the object itself) as well as different mass.

\begin{table}[H]
\centering
\small
\resizebox{\linewidth}{!}{%
\begin{tabular}{lcccc}
\toprule
Setting & Mass (kg) & Object $\mu$ & Table $\mu$ & Object spin \\
\midrule
Easy target & 5.0 & 0.5 & 0.5 & 0.0 \\
Hard target & $\mathcal{U}(5,10)$ & $\mathcal{U}(0.2,1.5)$ & $\mathcal{U}(0.2,1.5)$ & $\mathcal{U}(0,0.01)$ \\
\bottomrule
\end{tabular}
}
\caption{Target-domain evaluation settings used for the final success-rate comparison.}
\label{tab:evaluation-settings}
\end{table}
